\documentclass[letterpaper,11pt]{article}

\usepackage[letterpaper,left=0.80in,right=0.80in,top=0.60in,bottom=0.60in]{geometry}
\usepackage[hidelinks]{hyperref}
\usepackage[english]{babel}
\input{glyphtounicode}

\pagestyle{empty}
\urlstyle{same}
\raggedright
\setlength{\parindent}{0pt}
\linespread{1.0}
\pdfgentounicode=1

\begin{document}

\begin{center}
  {\LARGE \textbf{Aryan Miriyala}}\\[2pt]
  \small
  +1 419-315-0444 $|$ \href{mailto:aryanmiriyala@gmail.com}{aryanmiriyala@gmail.com} $|$ \href{https://www.linkedin.com/in/aryan-miriyala-7788a21b6/}{linkedin.com/in/aryan-miriyala}\\
  \href{https://github.com/aryanmiriyala}{github.com/aryanmiriyala}
\end{center}\vspace{8pt}

\today

\vspace{8pt}

Hiring Team\\
Otter.ai\\
Mountain View, CA

\vspace{8pt}

Dear Hiring Team,

\vspace{4pt}

I am applying for the Software Engineer, AI Agent \& LLM position at Otter.ai. With over 1 billion meetings transcribed and a mission to make conversations more valuable, Otter.ai sits at the intersection of conversational AI and practical productivity — exactly where I want to build.

\vspace{4pt}

In my current role at Actual Reality Technologies, I built production agent infrastructure from the ground up: a LiteLLM AI Gateway handling policy-controlled model routing with smart-router tiers for latency/cost optimization, guardrails, Langfuse tracing, Slack block alerts, CI release gates, and 44 red-team fixtures. I did this using AI coding tools daily — Claude Code, GitHub Copilot, Cursor, and Codex — to multiply velocity while maintaining rigorous review standards. I also built a remote MCP server enabling sandboxed execution for customer-facing AI tools with scoped security, and hardened authentication flows to eliminate redirect flashes in production AI workflows.

\vspace{4pt}

At SmartSolve Industries, I deployed role-specific AI agents for employee workflows and architected a full-stack onboarding tracker (Next.js, TypeScript, PostgreSQL, Drizzle ORM, SSO, auth middleware) from scratch. I applied Codex and Claude Code across implementation planning, debugging, and code review with human validation loops — the same rigorous, agent-assisted development approach your job description emphasizes.

\vspace{4pt}

On the AI quality and evaluation side, I developed Diff-Grounded PR Generation — an ICSME-published LLM evaluation pipeline that turns commits, diffs, and repository metadata into evidence-grounded PR descriptions, with automated evaluation for correctness, coverage, and clarity across 1,450 PRs. This directly mirrors your need for robust evaluation datasets, automated graders, regression tests, and release gates for AI quality.

\vspace{4pt}

For conversational knowledge systems, I built FalconGraph Search — a source-grounded RAG system that crawls resources, processes documents, generates embeddings, retrieves with FAISS vector search, and serves answers through a Next.js interface. This demonstrates end-to-end knowledge extraction, annotation, indexing, and conversational Q\&A — the core of Otter's conversational knowledge engine.

\vspace{4pt}

What draws me to Otter.ai is the opportunity to apply production AI agent infrastructure, LLM evaluation expertise, and conversational RAG experience at a scale where quality directly impacts millions of users. The focus on diagnosing failures across models, prompts, retrieval, and tools — and turning production traces into measurable improvements — matches how I operate: I build observable agent infrastructure, I define quality metrics for nondeterministic systems, and I use data to drive continuous improvement.

\vspace{4pt}

I would welcome the opportunity to discuss how my experience building AI Gateway infrastructure, LLM evaluation pipelines, and conversational AI agents can contribute to Otter AI Chat and the conversational knowledge engine.

\vspace{8pt}

Sincerely,

\vspace{4pt}

Aryan Miriyala

\end{document}